\section*{अध्याय \ref{ch:g12:metamorphic-rocks} --- कायांतरित चट्टानें}

\begin{solution}{exo:g12:metamorphic-rocks:1}
(a) \omterm{def:g12:metamorphic-rocks:polymorph}{बहुरूप}; (b) \omterm{def:g12:metamorphic-rocks:index-mineral}{समश्रेणी-रेखा}; (c) \omterm{def:g12:metamorphic-rocks:contact-regional}{परिमंडल}; (d) \omterm{def:g12:metamorphic-rocks:foliation}{पत्रण};
(e) \omterm{def:g12:metamorphic-rocks:lithostatic}{शैलस्थैतिक दाब}।
\end{solution}

\begin{solution}{exo:g12:metamorphic-rocks:2}
(1) संगमरमर; (2) क्वार्टज़ाइट; (3) स्लेट; (4) नाइस; (5) हॉर्नफ़ेल्स।
\end{solution}

\begin{solution}{exo:g12:metamorphic-rocks:3}
$P = \rho g z$: \qty{10}{km} पर, $2830 \times 9.81 \times \num{10000} \approx
\qty{2.78e8}{Pa} = \qty{0.278}{GPa}$; \qty{20}{km} पर, उसका दुगुना,
\qty{0.555}{GPa}।
\end{solution}

\begin{solution}{exo:g12:metamorphic-rocks:4}
(1) कायनाइट (\qty{400}{\celsius} पर कायनाइट--ऐंडालुसाइट रेखा लगभग \qty{0.25}{GPa}
के निकट है)। (2) ऐंडालुसाइट। (3) सिलिमेनाइट।
\end{solution}

\begin{solution}{exo:g12:metamorphic-rocks:5}
(1) स्लेट, फ़िलाइट, शिस्ट, नाइस। (2) शिस्ट में। (3) नाइस।
\end{solution}

\begin{solution}{exo:g12:metamorphic-rocks:6}
(1) $2830 \times 9.81 \times \num{35000} \approx \qty{9.72e8}{Pa}$, यानी
$\num{9.72e8}/101325 \approx \num{9600}$ वायुमंडल। (2) $1025 \times 9.81
\times 4000 \approx \qty{4.02e7}{Pa} = \qty{40.2}{MPa}$। (3) $972/40.2
\approx 24$ गुना।
\end{solution}

\begin{solution}{exo:g12:metamorphic-rocks:7}
(1) $z = \num{0.375e9}/(2830 \times 9.81) \approx \qty{13.5}{km}$।
(2) $(504 - 10)/13.5 \approx \qty{36.6}{\celsius}$ प्रति किलोमीटर। (3) नहीं:
चट्टान को आरेख के ठीक एक बिंदु पर बैठना पड़ता; ऐसी चट्टानें
दुर्लभ हैं।
\end{solution}

\begin{solution}{exo:g12:metamorphic-rocks:8}
(1) $23/3 \approx \qty{8}{mm}$। (2) लगभग 8 मिलियन वर्ष। (3) चट्टान
कई मिलियन वर्षों तक गहराई में, ऊँचे तापमान पर, रही।
\end{solution}

\begin{solution}{exo:g12:metamorphic-rocks:9}
(1) चार: गार्नेट, स्टॉरोलाइट, कायनाइट और सिलिमेनाइट \omterm{def:g12:metamorphic-rocks:index-mineral}{समश्रेणी-रेखाएँ}, इसी
क्रम में। (2) पश्चिम की ओर। (3) पश्चिमी सिरे पर, सिलिमेनाइट क्षेत्र में। उसकी
चट्टानों ने कायनाइट--सिलिमेनाइट रेखा पार की: लगभग \qty{20}{km}
(\qty{0.56}{GPa}) पर दबी होने के कारण वे लगभग \qty{590}{\celsius} से अधिक गरम थीं।
\end{solution}

\begin{solution}{exo:g12:metamorphic-rocks:10}
(1) पंक हॉर्नफ़ेल्स बन जाता है, चूना पत्थर संगमरमर, बलुआ पत्थर
क्वार्टज़ाइट। (2) वे बिना भींचे गरम हुईं: किसी दिष्ट दाब ने
उनके खनिजों को एक सीध में नहीं किया। (3) उसके चारों ओर की चट्टानें पहले से गरम हैं, इसलिए
मैग्मा उन्हें बस थोड़ा और गरम करता है।
\end{solution}

\begin{solution}{exo:g12:metamorphic-rocks:11}
(1) धक्के के आर-पार: उनके चपटे फलक पूर्व और पश्चिम की ओर मुख किए होते हैं, इसलिए शल्क
उत्तर--दक्षिण तलों में खड़े होते हैं। (2) ऊर्ध्वाधर, उत्तर--दक्षिण चलता हुआ।
(3) कैल्साइट और क्वार्ट्ज़ के कण न चपटे हैं, न लंबे, इसलिए भिंचाव में
पुनः क्रिस्टलित होने पर भी उनमें एक सीध में आने लायक कुछ नहीं है।
\end{solution}

\begin{solution}{exo:g12:metamorphic-rocks:12}
(1) $z = (T - 10)/30$, इसलिए $P = 0.0278\,(T - 10)/30 \approx
0.000927\,(T - 10)$। (2) $0.000927\,(T - 10) = 0.577 + 0.00209\,(T - 600)$
से $0.001163\,T \approx 0.668$ मिलता है, इसलिए $T \approx \qty{574}{\celsius}$, $P
\approx 0.000927 \times 564 \approx \qty{0.523}{GPa}$ और $z = (574 - 10)/30
\approx \qty{18.8}{km}$। (3) उस गहराई से ऊपर कायनाइट, उससे नीचे सिलिमेनाइट।
\end{solution}

\begin{solution}{exo:g12:metamorphic-rocks:13}
(1) $10 + 10 \times 35 = \qty{360}{\celsius}$। (2) $10 + 30 \times 35 =
\qty{1060}{\celsius}$। (3) \qty{35}{km} पर दाब लगभग
\qty{0.97}{GPa} है; \qty{360}{\celsius} पर वह बिंदु कायनाइट
क्षेत्र के भीतर गहरे पड़ता है। दूसरा बिंदु आरेख के बहुत दाईं ओर,
सिलिमेनाइट क्षेत्र में, उन तापमानों पर पड़ता है जहाँ भूपर्पटी की चट्टानें पिघलने लगती हैं।
(4) उन्हें निम्न तापमान पर ऊँचा दाब चाहिए, एक ऐसा संयोग जो केवल
\omterm{def:g11:subduction-zones:subduction}{क्षेपण क्षेत्र} में नीचे ले जाई गई ठंडी भूपर्पटी देती है।
\end{solution}

\begin{solution}{exo:g12:metamorphic-rocks:14}
(1) कायनाइट $1000/3.61 \approx \qty{277}{cm^3}$; सिलिमेनाइट $1000/3.24
\approx \qty{309}{cm^3}$; ऐंडालुसाइट $1000/3.15 \approx \qty{317}{cm^3}$।
(2) $(317 - 277)/317 \approx 0.127$, लगभग \qty{13}{\percent}। (3) ऊँचे
दाब पर परमाणुओं की सबसे सघन व्यवस्था का पक्ष लिया जाता है; कायनाइट
उन्हीं परमाणुओं को सबसे कम जगह में समेटता है।
\end{solution}

\begin{solution}{exo:g12:metamorphic-rocks:15}
(1) \qty{20}{km} पर \qty{0.555}{GPa}, \qty{25}{km} पर \qty{0.694}{GPa}।
(2) $T = 600 + (P - 0.577)/0.00209$: \qty{0.555}{GPa} पर लगभग
\qty{590}{\celsius} और \qty{0.694}{GPa} पर \qty{656}{\celsius}।
\qty{700}{\celsius} पर चट्टानें दोनों से अधिक गरम थीं, सिलिमेनाइट
क्षेत्र में: संगत। (3) $(700 - 10)/20 \approx \qty{34.5}{\celsius}$ और
$(700 - 10)/25 \approx \qty{27.6}{\celsius}$ प्रति किलोमीटर, सामान्य
प्रवणता के निकट। (4) 40 मिलियन वर्षों में $\qty{20}{km}$ से $\qty{25}{km}$:
प्रति मिलियन वर्ष 0.50 से \qty{0.63}{km}, यानी प्रति वर्ष 0.50 से \qty{0.63}{mm}।
\end{solution}

\begin{solution}{pb:g12:metamorphic-rocks:1}
\textbf{1.} वह पत्रित है, और उसके अभ्रक इतने बड़े हैं कि दिखाई दें।
अभ्रक, मस्कोवाइट और बायोटाइट, \omterm{def:g12:metamorphic-rocks:foliation}{शिस्टाभता} बनाते हैं।

\textbf{2.} \ce{Al2SiO5}; वे \omterm{def:g12:metamorphic-rocks:polymorph}{बहुरूप} हैं।

\textbf{3.} एक किलोग्राम के आयतन $3.15/3.61 \approx
0.873$ के अनुपात में हैं: कायनाइट लगभग \qty{13}{\percent} छोटा है। दाब
कायनाइट के पक्ष में होता है।

\textbf{4.} कायनाइट क्षेत्र। नीचे जाते हुए उसने बायोटाइट,
गार्नेट, स्टॉरोलाइट और कायनाइट \omterm{def:g12:metamorphic-rocks:index-mineral}{समश्रेणी-रेखाएँ} पार कीं।

\textbf{5.} लगभग 8 मिलियन वर्ष।

\textbf{6.} $P = 0.577 + 0.00209 \times 0 = \qty{0.577}{GPa}$।

\textbf{7.} \qty{0.70}{GPa}, \qty{0.577}{GPa} से ऊपर है, कायनाइट--सिलिमेनाइट
रेखा की कायनाइट वाली ओर; वह ऐंडालुसाइट क्षेत्र से भी बहुत ऊपर है,
जो \qty{600}{\celsius} पर लगभग \qty{0.23}{GPa} से नीचे है।

\textbf{8.} \qty{0.577}{GPa}।

\textbf{9.} $0.70 = 0.577 + 0.00209\,(T - 600)$, इसलिए $T = 600 +
0.123/0.00209 \approx \qty{659}{\celsius}$।

\textbf{10.} \qty{7.0}{kbar}; $\num{7.0e8}/101325 \approx \num{6900}$
वायुमंडल।

\textbf{11.} $\rho g = 2830 \times 9.81 \approx \qty{27800}{Pa}$ प्रति मीटर,
\qty{27.8}{MPa} प्रति किलोमीटर।

\textbf{12.} $z = \num{0.70e9}/\num{27762} \approx \qty{25200}{m}$, लगभग
\qty{25.2}{km}।

\textbf{13.} $\num{0.577e9}/\num{27762} \approx \qty{20.8}{km}$।

\textbf{14.} $\num{0.70e9}/(2700 \times 9.81) \approx \qty{26.4}{km}$, लगभग
\qty{5}{\percent} अधिक गहरा ($26.4/25.2 \approx 1.048$)।

\textbf{15.} भूपर्पटी में: \qty{25}{km} भूपर्पटी की सामान्य मोटाई से भी
कम है।

\textbf{16.} $(600 - 10)/25.2 \approx \qty{23.4}{\celsius}$ प्रति किलोमीटर।

\textbf{17.} यह किसी \omterm{def:g11:subduction-zones:subduction}{क्षेपण क्षेत्र} (प्रति किलोमीटर
\qty{10}{\celsius} से कम) से बहुत गरम है और किसी सामान्य भूपर्पटी
(\qty{30}{\celsius}) से थोड़ा ठंडा, किसी ज्वालामुखीय क्षेत्र (40 से \qty{50}{\celsius}) से बहुत दूर:
किसी पर्वत-पट्टी की मोटी हुई भूपर्पटी में गहरे दबी चट्टान, \omterm{def:g12:metamorphic-rocks:contact-regional}{क्षेत्रीय
कायांतरण} द्वारा।

\textbf{18.} $10 + 30 \times 25.2 \approx \qty{766}{\celsius}$। आगे बढ़ाई गई
कायनाइट--सिलिमेनाइट रेखा उस तापमान पर $0.577 + 0.00209 \times 166
\approx \qty{0.92}{GPa}$ पर है; \qty{0.70}{GPa} उससे नीचे,
सिलिमेनाइट क्षेत्र में है। सामान्य \omterm{def:g11:lithosphere-asthenosphere:geotherm}{भूतापीय वक्र} सिलिमेनाइट बनाता।

\textbf{19.} $\qty{25.2}{km} = \qty{2.52e7}{mm}$; प्रति वर्ष \qty{0.5}{mm} पर,
लगभग 50 मिलियन वर्ष।

\textbf{20.} ऊपर उठते हुए चट्टान ठंडी हुई, इसलिए उसकी अभिक्रियाएँ धीमी हो गईं, और
उलटी अभिक्रियाओं को जिस जल की आवश्यकता होती है, वह नीचे जाते समय ही
निकाल दिया गया था।

\textbf{21.} कायनाइट शिस्ट लगभग \textbf{\qty{25}{km}
नीचे}, लगभग \qty{600}{\celsius} पर, क्रिस्टलित हुआ।
\end{solution}
